% sections/03_mathematical_framework.tex
\section{Mathematical Framework}\label{sec:math-framework}

This section defines the core formulae that transform raw bilateral flow data
into normalised concentration scores.  All notation introduced here applies
throughout the remainder of the document.

\subsection{Notation}\label{sec:math:notation}

Let \(i\) index countries (\(i = 1, \ldots, N\); \(N = 27\) for v\,1.0).
Let \(a\) index thematic axes (\(a = 1, \ldots, 6\)).  Where an axis defines
two channels, they are denoted \(\mathrm{ch} \in \{A, B\}\).  Not all axes
use a two-channel decomposition; axis-specific structures are detailed in
\cref{sec:axis-specifications}.

For a given country~\(i\), axis~\(a\), and channel~\(\mathrm{ch}\), the
partner-share vector is
\(\mathbf{s}_i^{(\mathrm{ch})} = (s_{i,1}^{(\mathrm{ch})}, \ldots,
s_{i,J}^{(\mathrm{ch})})\), where \(s_{i,j}^{(\mathrm{ch})}\) is the share
of partner~\(j\) in the total flow and
\(\sum_{j} s_{i,j}^{(\mathrm{ch})} = 1\).

\subsection{Concentration Formula}\label{sec:math:hhi}

Concentration is measured by the Herfindahl--Hirschman Index applied to the
partner-share vector.  For country~\(i\) on channel~\(\mathrm{ch}\):
%
\begin{equation}\label{eq:hhi}
  C_i^{(\mathrm{ch})}
  \;=\;
  \sum_{j=1}^{J} \bigl(s_{i,j}^{(\mathrm{ch})}\bigr)^{\!2}
\end{equation}
%
The resulting value lies on the interval \((0, 1]\).  A value of~\(1\)
indicates total dependence on a single partner; values approaching~\(0\)
indicate dispersion across many partners.

\paragraph{Relation to the classical \HHI.}
The traditional \HHI in industrial-organisation literature is computed on
percentage shares and ranges from~\(0\) to~\num{10000}.  The \ISI operates
on fractional shares, yielding a normalised \HHI on \((0, 1]\).  The two
forms are related by a factor of~\num{10000}.

\subsection{Cross-Channel Aggregation}\label{sec:math:cross-channel}

For axes that define two channels, the axis-level score is the arithmetic
mean of the channel-level scores:
%
\begin{equation}\label{eq:dual-channel}
  M_i^{(a)}
  \;=\;
  0.5\, C_i^{(A)} \;+\; 0.5\, C_i^{(B)}
\end{equation}
%
Equal weighting is a deliberate design choice (Principle~\textbf{P4}).  No
empirical basis exists for assigning asymmetric importance to supply versus
destination concentration in a general-purpose index.

Where only one channel is computable (fallback basis \texttt{A\_ONLY} or
\texttt{B\_ONLY}), the axis score equals the surviving channel score.

\subsection{Category-Weighted Channel~B}\label{sec:math:weighted-b}

For axes where Channel~\(B\) encompasses multiple product sub-categories or
capability blocks (see \cref{sec:axis-specifications}), a category-weighted
variant is used.  Let \(k = 1, \ldots, K\) index the sub-categories and let
\(V_i^{(k)}\) denote the total flow value of sub-category~\(k\) for
country~\(i\).  Then:
%
\begin{equation}\label{eq:weighted-b}
  C_i^{(B)}
  \;=\;
  \frac{\displaystyle\sum_{k=1}^{K} C_i^{(B,k)}\, V_i^{(k)}}
       {\displaystyle\sum_{k=1}^{K} V_i^{(k)}}
\end{equation}
%
where \(C_i^{(B,k)}\) is the \HHI of sub-category~\(k\).  This ensures that
small-volume sub-categories do not dominate the axis score.

\subsection{Single-Channel Fuel Averaging}\label{sec:math:fuel-avg}

Axis~2 (Energy) does not use a two-channel decomposition.  Instead, a
per-fuel \HHI is computed for each fossil-fuel category and the axis score is
the arithmetic average across fuels:
%
\begin{equation}\label{eq:fuel-avg}
  M_i^{(2)}
  \;=\;
  \frac{1}{|\mathcal{F}_i|}
  \sum_{f \in \mathcal{F}_i} C_i^{(f)}
\end{equation}
%
where \(\mathcal{F}_i\) is the set of fuel categories for which
country~\(i\) has non-zero imports (see \cref{sec:axis:energy}).

\subsection{Composite Score}\label{sec:math:composite}

The composite \ISI score for country~\(i\) is the unweighted arithmetic mean
of the six axis-level scores:
%
\begin{equation}\label{eq:composite}
  \mathrm{ISI}_i
  \;=\;
  \frac{A_1 + A_2 + A_3 + A_4 + A_5 + A_6}{6}
\end{equation}
%
where \(A_a \equiv M_i^{(a)}\) for notational compactness.  The composite
inherits the unit-interval range of the axis scores.

\paragraph{Equal weighting.}
No differential axis weights are applied.  Differential weighting would
require an empirical or normative justification for asserting that one domain
matters more than another for \emph{all} countries at \emph{all} times.  No
credible basis for such an assertion exists.  The equal-weight rule is
transparent, reproducible, and neutral across policy domains.

\subsection{Rounding}\label{sec:math:rounding}

All published scores---axis-level and composite---are rounded to eight
decimal places using the \emph{round-half-to-even} convention (IEEE~754
banker's rounding).  The rounding precision is set by the backend constant
\texttt{ROUND\_PRECISION\,=\,8}.  Rounding is applied \emph{before}
classification and \emph{before} hashing.  Intermediate calculations retain
full floating-point precision.

\subsection{Classification Thresholds}\label{sec:math:thresholds}

The rounded composite score is mapped to a four-tier ordinal classification:

\begin{table}[htbp]
  \centering\small
  \caption{Classification thresholds.}\label{tab:thresholds}
  \begin{tabular}{@{}l l@{}}
    \toprule
    \textbf{Composite range}     & \textbf{Classification} \\
    \midrule
    \(\mathrm{ISI}_i \geq 0.50\) & \texttt{highly\_concentrated}     \\
    \(0.25 \leq \mathrm{ISI}_i < 0.50\) & \texttt{moderately\_concentrated} \\
    \(0.15 \leq \mathrm{ISI}_i < 0.25\) & \texttt{mildly\_concentrated}     \\
    \(\mathrm{ISI}_i < 0.15\)    & \texttt{unconcentrated}          \\
    \bottomrule
  \end{tabular}
\end{table}

The threshold values are adapted from the US Department of Justice \HHI
guidelines for market concentration \parencite{doj2023merger}, rescaled to
the unit interval.  They are defined in \texttt{registry.json} and enforced
by a single classification function in \texttt{methodology.py}.  The
thresholds are fixed parameters of \ISI v\,1.0.  Any modification to the
threshold values constitutes a major-version increment under the governance
protocol (\cref{sec:governance}).